%%%%%%%%%%%%%%%%%%%%%%%%%%%%%%%%%
%
%       EXERCÍCIO
%
%%%%%%%%%%%%%%%%%%%%%%%%%%%%%%%%%

\ifdefstring{\atividade}{prova}{%
    \ifdefstring{\modo}{objetivo}{%
        \renewcommand{\valorquestao}{\ValorQObj\ ponto}
    }{%
        \renewcommand{\valorquestao}{\ValorQDisc\ pontos}
    }%
}%

\begin{exercicioBanco}[\valorquestao]
Determine o gráfico da função afim \(f: \bbR \to \bbR\), dada por \(f(x) = 3x + 2\).

% Define as alternativas
\newcommand{\alternativas}{%
\begin{center}
    \begin{tabular}{C{20pt}C{140pt}C{20pt}C{140pt}C{20pt}C{140pt}}
        (a) 
        &
        \begin{tikzpicture}[scale=0.5]
            % Eixos
            \draw[->] (-3.0,0) -- (3.0,0) node[right] {\(x\)};
            \draw[->] (0,-3.0) -- (0,3.0) node[above] {\(y\)};

            % Gráfico da função f(x) = |x+1| - 2
            \draw[domain=-2.8:2.8, thick, blue] plot (\x, {\x});
        \end{tikzpicture}
        &
        (b) 
        &
        \begin{tikzpicture}[scale=0.5]
            % Eixos
            \draw[->] (-3.0,0) -- (3.0,0) node[right] {\(x\)};
            \draw[->] (0,-3.0) -- (0,3.0) node[above] {\(y\)};

            % Gráfico da função
            \draw[domain=-1.5:0.5, thick, blue] plot (\x, {3*(\x) + 2});

            % Marcações da raiz
            \node[above left] at (-0.67,0) {\small\(-\frac{2}{3}\)};
            \node at (-0.67,0) {\small\(\bullet\)};

            % Marcação do ponto b
            \node[right] at (0,2) {\small\(2\)};
            \node at (0,2) {\small\(\bullet\)};
        \end{tikzpicture}
        &
        (c)
        &
        \begin{tikzpicture}[scale=0.5]
            % Eixos
            \draw[->] (-3.0,0) -- (3.0,0) node[right] {\(x\)};
            \draw[->] (0,-3.0) -- (0,3.0) node[above] {\(y\)};

            % Gráfico da função
            \draw[domain=-0.5:1.5, thick, blue] plot (\x, {3*(\x) - 2});

            % Marcações da raiz
            \node[below right] at (0.67,0) {\small\(\frac{2}{3}\)};
            \node at (0.67,0) {\small\(\bullet\)};

            % Marcação do ponto b
            \node[left] at (0,-2) {\small\(-2\)};
            \node at (0,-2) {\small\(\bullet\)};
        \end{tikzpicture}
        \\
        (d)
        &
        \begin{tikzpicture}[scale=0.5]
            % Eixos
            \draw[->] (-3.0,0) -- (3.0,0) node[right] {\(x\)};
            \draw[->] (0,-3.0) -- (0,3.0) node[above] {\(y\)};

            % Gráfico da função
            \draw[domain=-2.7:1.5, thick, blue] plot (\x, {\x + 2});

            % Marcações da raiz
            \node[above left] at (-2,0) {\small\(-2\)};
            \node at (-2,0) {\small\(\bullet\)};

            % Marcação do ponto b
            \node[right] at (0,2) {\small\(2\)};
            \node at (0,2) {\small\(\bullet\)};
        \end{tikzpicture}
        &
        (e) 
        &
        \begin{tikzpicture}[scale=0.5]
            % Eixos
            \draw[->] (-3.0,0) -- (3.0,0) node[right] {\(x\)};
            \draw[->] (0,-3.0) -- (0,3.0) node[above] {\(y\)};

            % Gráfico da função
            \draw[domain=-0.5:1.5, thick, blue] plot (\x, {-3*(\x) + 2});

            % Marcações da raiz
            \node[above right] at (0.67,0) {\small\(\frac{2}{3}\)};
            \node at (0.67,0) {\small\(\bullet\)};

            % Marcação do ponto b
            \node[left] at (0,2) {\small\(2\)};
            \node at (0,2) {\small\(\bullet\)};
        \end{tikzpicture}
        &&
    \end{tabular}
\end{center}
}

% Define a resposta correta
\newcommand{\resposta}{B}

% Lógica condicional para exibição
\ifdefstring{\atividade}{lista}{%
    \alternativas
    \vspace{0.5em}
    
    \noindent\textbf{Resposta:} letra \textbf{\resposta}.
}{%
    \ifdefstring{\modo}{objetivo}{%
        \alternativas
    }{%
        % modo = discursiva → não mostra alternativas
    }
}
\end{exercicioBanco}

